\section{Conclusion}
\label{sec:conclusion}

We set out to make retrieval-augmented program repair work better by making
retrieval more relevant, and instead produced a cautionary decomposition of
why that framing fails. Structured failure-aware reranking can improve a
retrieval-relevance proxy, but the gain is fragile across benchmarks, and
even where it holds it does not convert to repair: downstream effects are
neutral-to-slightly-negative across six models and four benchmarks, bounded
within $\pm 10\pp$. A planted-corpus experiment locates the binding
constraint: when a sufficiently relevant example exists in the corpus, plain
dense retrieval already surfaces it and repair improves by double digits;
when it does not --- the common case --- no amount of reranking
sophistication helps, and our reranker even selects the perfect example
slightly less often than dense retrieval. Coverage dominates ranking. On real
repository bugs the same pattern holds directionally, with a powered
replication in progress.

For practitioners, the guidance is concrete: before investing in retrieval
scoring for repair, invest in the corpus --- coverage and curation are where
the recoverable headroom sits, and the headroom scales with how much the
model currently fails. For evaluators, our results argue for
contamination-controlled corpora, post-processing ablations, and equivalence
bounds as defaults in RAG-repair studies. All artifacts, including the
planted-corpus builders and the CI harness for powered real-bug replication,
are released for reuse.
